\section{Semantics of Federated SPARQL Queries}\label{sec:semantic_federated_q}

Before investigating query containment, we must establish under which semantics the different queries we study are evaluated.
This choice determines whether and how duplicate results can arise, which directly impacts how query containment is resolved.

\subsection{Single Endpoint Execution}\label{sec:single_endpoint}
SPARQL queries on single sources are evaluated under bag-set semantics.
The underlying \ac{KG} operates under set semantics~\cite{w3ConceptsAbstract}, whereas SPARQL operators produce bags.
The choice between bag and set semantics has a significant impact on the complexity and decidability of query containment~\cite{Chaudhuri1993}.
For relational database settings, \citeauthor{Chaudhuri1993}~\cite{Chaudhuri1993, Afrati2010} established that for \acp{CQ} only projection can increase the multiplicity of a result.
The same property holds for SPARQL \acp{CQ}, where the only body operator, \texttt{JOIN}, cannot raise the multiplicity, as in the relational case~\cite{Chaudhuri1993, Afrati2010}; projection at the head therefore remains the only operator that can do so.
In the context of \ac{UCQ}, the multiplicity can additionally be increased by \texttt{UNION} and by the combination of \texttt{UNION} and \texttt{JOIN}, as defined in \Cref{def_union_multiplicity,def_join_multiplicity}.
\ac{UCQ} is not the subject of our study; however, in the next sections we introduce \ac{UCFQ}, a subclass of \ac{UCQ}, together with the semantics of federated queries.



\subsection{Federated Queries}
A \ac{KG} is stored under set semantics, and querying a single endpoint operates under bag-set semantics.
A federated query, however, combines solution mappings drawn from several \acp{KG} at once, so it is not obvious whether spanning multiple members can introduce duplicate results.
We therefore examine the semantics of federated query evaluation.
The SPARQL federated query specification~\cite{w3SPARQLFederated} defines the execution of federated queries with \texttt{SERVICE} clauses as a join between the bag of solution mappings from the preceding execution and the bag of solution mappings returned by the SPARQL service.
The \ac{GP} inside a \texttt{SERVICE} clause has no projection, so each federation member returns a set of solution mappings.
A \texttt{SERVICE} clause is nested in the query through a \texttt{JOIN}, which cannot increase the multiplicity.
Hence, even when some federation members hold duplicate information, the federated query produces no duplicate results, and \ac{CQ} evaluation with \texttt{SERVICE} clauses remains under bag-set semantics.

\subsection{\acl{UCFQ}}\label{sec:ucfq}

We define a particular class of federated queries, useful for capturing duplicate information held by different federation members.
Each triple pattern of the query body is evaluated across several members and the partial results are joined.

\begin{definition}[\acl{UCFQ}]\label{def_ucfq}
A \ac{UCFQ} is a multiway join $\fmj(U_1, \dots, U_n)$ where each operand $U_i = \fmu(\freq_{I_1}^{\mathit{tp}_i}, \dots, \freq_{I_{k_i}}^{\mathit{tp}_i})$ is a multiway union evaluating the same triple pattern $\mathit{tp}_i$, up to a renaming of the variables local to $U_i$, at a subset of federation members $I_1, \dots, I_{k_i}$, optionally followed by a projection.
\end{definition}

A variable is \emph{local} to $U_i$ when the head discards it and no other operand of the join binds it.
Renaming the local variables of a branch, distinct variables to distinct variables, maps its solution mappings one to one and preserves their multiplicity, and no other operator of the query binds them, so the result bag is unchanged.
Branches naming them differently therefore evaluate the same triple pattern, which matters as a query engine binding a variable at a single branch may name it freely.

A \ac{CQ} is the special case where every union ranges over a single member, and every \ac{UCFQ} is a \ac{UCQ}, so the \acp{UCFQ} lie strictly between the \acp{CQ} and the \acp{UCQ}.
In particular, each triple pattern of a \ac{CQ} is the single-member union $\fmu(\freq_{\mathit{fm}}^{\mathit{tp}})$ over the member $\mathit{fm}$ at which it is evaluated.

\acp{UCFQ} arise under automatic source selection, in particular \ac{ESA} (\Cref{sec:related_works_fed}).
Consider a federation of two members, a job registry and a social network, that redundantly store the job and birth date of a person.
The projection-free \ac{CQ} retrieving each person's job and birth date (\Cref{fig:ucfq-example}, top) would return a set on a single endpoint, but under \ac{ESA} it is evaluated as the \ac{UCFQ} (\Cref{fig:ucfq-example}, bottom), unioning each triple pattern across both members before joining the results.
Since both members hold the same solution for a person, the join returns that person with multiplicity four, even though every member stores a set --- a duplication that does not arise with explicit \texttt{SERVICE} clauses in a \ac{CQ}.%
\footnote{An executable example of the semantics of federated queries, evaluated with the Comunica~\cite{taelman_iswc_resources_comunica_2018} query engine, is available at \url{https://github.com/constraintAutomaton/demo-federated-query-semantic}.}

\begin{figure}
  \textbf{The \ac{CQ} as written}
  \lstinputlisting[style=sparql]{code/ucfq_cq.rq}
  \textbf{The executed query under \ac{ESA}}
  \lstinputlisting[style=sparql]{code/ucfq_example.rq}
  \caption{A \ac{CQ} evaluated as a \ac{UCFQ} under \ac{ESA}.}
  \label{fig:ucfq-example}
\end{figure}

\ac{ESA} is the particular case of \ac{UCFQ} in which every union branch ranges over the entire federation $\{\mathit{fm}_1, \dots, \mathit{fm}_m\}$, formalized as follows:
\begin{equation}
  \fmj\bigl(\fmu(\freq_{\mathit{fm}_1}^{\mathit{tp}_1}, \dots, \freq_{\mathit{fm}_m}^{\mathit{tp}_1}), \dots, \fmu(\freq_{\mathit{fm}_1}^{\mathit{tp}_n}, \dots, \freq_{\mathit{fm}_m}^{\mathit{tp}_n})\bigr)
  \label{eq_esa}
\end{equation}

Most federated engines produce plans of this same join-over-union shape---for instance FedX~\cite{schwarte2011fedx}, SemaGrow~\cite{charalambidis2015semagrow}, and CostFed~\cite{saleem2018costfed}---which can likewise be rewritten as \acp{UCFQ}.
Automatic source selection can, however, produce \ac{CQ} plans of arbitrary shape --- for instance, union-over-join plans~\cite{aimonier2024fedup} --- and \ac{UCFQ} therefore does not capture every \ac{CQ} under automatic source selection.

Containment of \acp{UCQ} under bag-set semantics is undecidable (\Cref{sec:related_works_qc}), but this does not settle the status of the \acp{UCFQ}, which form a strict subclass.
The next section introduces the \ac{BFR}, a transformation showing that evaluating a \ac{UCFQ} is equivalent to evaluating a \ac{CQ} over one \ac{BKG} per union branch under bag semantics.



\subsection{\acl{BFR}}\label{sec_bag_federation_reduction}

We propose the \ac{BFR}, a transformation that rewrites a \ac{UCFQ} as a \ac{CQ} evaluated over one \ac{BKG} per union branch (\Cref{def_bkg}).
Intuitively, a virtual operator aggregates the members of each branch's sub-federation into a \ac{BKG} and evaluates the branch's triple pattern over it.

\subfile{proofs/bag_federation_reduction}

\subsection{A Bag-Semantic Reading of FedQPL}\label{sec_fedqpl}

\Cref{sec:single_endpoint,sec:ucfq} established that SPARQL evaluates under bag-set semantics: the underlying \acp{KG} are sets, while the operators combining their results return bags, so that unioning a triple pattern across federation members duplicates solution mappings (\Cref{def_union_multiplicity}).
We expressed the queries this produces in the notation of FedQPL (\Cref{def_ucfq}), which carries no such semantics: \citeauthor{Cheng2021} evaluate every operator over sets of solution mappings, and $\fmu$ denotes set union~\cite[Def.~6]{Cheng2021}.
We read $\fmu$ and $\fmj$ as their bag counterparts (\Cref{def_fedqpl}), which aligns them with the SPARQL operators they lift and places FedQPL under the bag-set semantics of \Cref{sec:single_endpoint}.
Other readings are possible, and we return to them below.
\citeauthor{Cheng2021} give a bag semantics as future work~\cite{Cheng2021}, and to the best of our knowledge no such extension has been published; we set out here what one has to settle.

The federation semantics of FedQPL evaluates a \ac{GP} over the set union $G_{\mathit{union}}$ of the \acp{KG} of the federation members~\cite[Def.~3]{Cheng2021}, following the stated intention that ``the result of a SPARQL query over such a federation should be the same as if the query was executed over the union of all the RDF data available in all the federation members''.
Correctness of a source assignment is defined against that graph~\cite[Def.~7]{Cheng2021}.
Under set semantics the exhaustive source assignment is correct~\cite[Prop.~1]{Cheng2021}; under bag-set semantics it is not.

\begin{proposition}\label{prop_fedqpl_esa_incorrect}
  Read under bag-set semantics, the exhaustive source assignment~\cite[Def.~9]{Cheng2021} is not correct in the sense of~\cite[Def.~7]{Cheng2021}.
\end{proposition}

\begin{proof}
  $G_{\mathit{union}}$ is a set, so a \ac{BGP} returns each of its solution mappings once over it (\Cref{sec:single_endpoint}).
  The \ac{ESA} expression of \Cref{eq_esa}, on the other hand, returns $\mu$ with multiplicity $\prod_{i} |\{\mathit{fm} \in F : \mu(\mathit{tp}_i) \in G_{\mathit{fm}}\}|$ by \Cref{def_union_multiplicity,def_join_multiplicity}.
  The federation of \Cref{fig:ucfq-example} makes this product four, whereas $\eval{B}{G_{\mathit{union}}}$ returns that solution mapping once.
\end{proof}

There are two ways to preserve~\cite[Def.~3]{Cheng2021} as stated.
The first assumes that no triple is duplicated across the federation members, so that the set union of their \acp{KG} coincides with the bag union.
Such an assumption is hard to verify in the open-Web setting and cannot always be enforced.
The second introduces a federation-level union operator that deduplicates the contributions of distinct members.
That operator is not a set union, but a different and more complex type of union.
A request may itself return a bag, since the request language of a SPARQL endpoint is the full graph pattern language, and those multiplicities must survive.
The operator must therefore deduplicate across members while leaving the multiplicities of each member's own response untouched, which requires the algebra to track the member every solution mapping comes from --- an operator SPARQL does not have, and whose addition changes the algebra rather than extending it.
We adopt neither.

We propose instead to use a \ac{BKG} (\Cref{def_bkg}) in place of the set union.
The federation is then a bag rather than a set, so the reading moves from bag-set to bag semantics, which is the transition the \ac{BFR} performs (\Cref{sec_bag_federation_reduction}).
\Cref{prop_duplicate_fed_mem}, taken with $F_i = F$ for every branch, restores the correctness of the exhaustive source assignment, and the intention of \citeauthor{Cheng2021} survives with the set union of the \acp{KG} replaced by their bag union.
This reading keeps the \texttt{UNION} of SPARQL and adds no operator to the algebra.
It also leaves the duplication of a triple across federation members observable in the result bag, so that a query can determine how many members hold a given triple --- information that serves domain purposes as well as the checking of integrity constraints over a federation.
We adopt it for the remainder of the paper.
